\bookchapter{Kapitel 54: Projekt 1 – Mein Haushaltsbuch (Lokalversion)}

\chapterintro{Willkommen bei deinem ersten “richtigen Projekt”. Dieses Haushaltsbuch geht zwar nicht ins Netz, aber es hat alles, was dazugehört: eine Buchung erfassen, den Verlauf ansehen, monatliche Auswertungen, Löschen und Ändern. Wenn du fertig bist, hast du den kompletten Ablauf “Daten Logik Oberfläche” einmal durchlaufen. Und: Nicht sofort den Code abschreiben – erst klar denken, dann anpacken. Das ist unser Prinzip von Anfang an.}

\booksection{Über die Anforderungen nachdenken: erst fünf Fragen stellen}

Bevor du loslegst, sortiere deine Idee mit den “fünf Anforderungsfragen”:

\begin{center}
\begin{tabularx}{\linewidth}{|X|X|}
\hline
\textbf{Frage} & \textbf{Antwort} \tabularnewline
\hline
Who (für wen?) & Für dich selbst, auch für Familie und Freunde \tabularnewline
What (was?) & Einnahmen und Ausgaben erfassen, Verlauf ansehen, monatliche Summen, Auswertung nach Kategorien \tabularnewline
Why (warum?) & Übung + herausfinden, wohin das Geld fließt \tabularnewline
How (wie?) & Kotlin + Compose + Room + MVVM, rein lokal ohne Backend \tabularnewline
Constraint (Rahmenbedingung) & In einer Woche fertig; mindestens Android 8 unterstützen; auch ohne Store-Veröffentlichung nach Store-Maßstab schreiben \tabularnewline
\hline
\end{tabularx}
\end{center}

\begin{center}\small\color{gray}Die fünf Anforderungsfragen des Haushaltsbuchs\end{center}

Dann zerlegst du es in kleine Schritte, von denen jeder einzeln “benutzbar” ist: (1) Room-Datenschicht bauen (eine Buchung speichern können); (2) Listenseite (Verlauf ansehen können); (3) Erfassungsseite (Buchung hinzufügen können); (4) Löschen und Bearbeiten; (5) Monats-Karte. Nach jedem Schritt ist die App ein “brauchbares Ding”.

\booksection{Die Datenschicht: Den “Buchungen” ein Zuhause geben}

Definiere zuerst, wie eine Buchung aussieht – die Felder sollen reichen, aber nicht gierig sein:

\begin{lstlisting}[language=Kotlin]
@Entity(tableName = "records")
data class Record(
    @PrimaryKey(autoGenerate = true) val id: Long = 0,
    val type: String,        // "income" Einnahme | "expense" Ausgabe
    val amount: Double,      // Betrag (Yuan)
    val category: String,    // Kategorie: Essen, Verkehr, Gehalt …
    val note: String = "",   // Notiz (optional)
    val date: String         // "2026-08-15"
)
\end{lstlisting}

In den Dao kommen zwei “Statistik-Abfragen” – sie sind die Seele des Haushaltsbuchs:

\begin{lstlisting}[language=Kotlin]
@Dao
interface RecordDao {
    @Query("SELECT * FROM records ORDER BY date DESC, id DESC")
    fun getAll(): Flow<List<Record>>

    @Insert
    suspend fun insert(record: Record)

    @Delete
    suspend fun delete(record: Record)

    @Query("SELECT COALESCE(SUM(amount),0) FROM records WHERE type='income' AND date LIKE :month || '%'")
    suspend fun monthIncome(month: String): Double   // month hat das Format "2026-08"

    @Query("SELECT COALESCE(SUM(amount),0) FROM records WHERE type='expense' AND date LIKE :month || '%'")
    suspend fun monthExpense(month: String): Double
}
\end{lstlisting}

\begin{tcbinfo}{Tipp: Warum Double für Beträge?}
Seriöse Finanz-Apps nutzen BigDecimal oder “Int in Cent”, um Rundungsfehler zu vermeiden (das alte Problem 0,1+0,2$\neq$0,3). Für ein Übungsprojekt reicht Double völlig; baust du eines Tages ein echtes Finanzprodukt, rechne in Cent – eine Falle, in die unzählige Vorgänger getreten sind.

\end{tcbinfo}

\booksection{MVVM: Oberfläche und Logik ziehen getrennte Wege}

Die Oberfläche direkt auf die Datenbank zugreifen lassen? Läuft zwar – aber der Oberflächen-Code wird lang und furchtbar. Die MVVM-Architektur lässt jeden seine Rolle spielen: View (Oberfläche) zeichnet nur, ViewModel (Logikzentrale) verwaltet die Daten, Model (Room) kümmert sich um die Speicherung.

View (Compose-Oberfläche)

ViewModel (Zustand und Logik)

Model (Datenlayer mit Room)

\begin{center}\small\color{gray}Abbildung 5: Die drei MVVM-Schichten – Daten fließen in eine Richtung, keine Schicht greift über\end{center}

\begin{lstlisting}[language=Kotlin]
class RecordViewModel(private val dao: RecordDao) : ViewModel() {

    // wandelt den Flow in einen von der UI beobachtbaren Zustand
    val records = dao.getAll()
        .stateIn(viewModelScope, SharingStarted.WhileSubscribed(), emptyList())

    // Nutzeraktionen: suspend-Funktionen aufrufen
    fun addRecord(type: String, amount: Double, category: String, note: String, date: String) {
        viewModelScope.launch {
            dao.insert(Record(type = type, amount = amount,
                category = category, note = note, date = date))
        }
    }

    fun deleteRecord(record: Record) {
        viewModelScope.launch { dao.delete(record) }
    }
}
\end{lstlisting}

In der Oberfläche (Compose) so: val vm: RecordViewModel = viewModel(), dann val list by vm.records.collectAsState() – sobald sich die Daten ändern, aktualisiert sich die Liste von selbst, du musst nichts tun. Das ist die geballte Kraft von MVVM + Flow + Compose.

\booksection{Oberfläche: Listenseite + Erfassungsseite}

Für die Listenseite nimm LazyColumn (träge ladende Liste) – selbst bei riesigen Datenmengen ruckelt nichts:

\begin{lstlisting}[language=Kotlin]
@Composable
fun RecordListScreen(vm: RecordViewModel) {
    val records by vm.records.collectAsState()

    Column {
        // oben die Zusammenfassungs-Karte (Einnahmen - Ausgaben = Saldo)
        SummaryCard(records)

        // Liste der Buchungen
        LazyColumn {
            items(records, key = { it.id }) { record ->
                RecordRow(record) { vm.deleteRecord(record) }
            }
        }

        // schwebender Button "Ausgabe erfassen"  öffnet den Eingabedialog
        FloatingActionButton(onClick = { showAddDialog = true }) {
            Text("+")
        }
    }
}

@Composable
fun RecordRow(record: Record, onDelete: () -> Unit) {
    Row(Modifier.fillMaxWidth().padding(12.dp)) {
        Text(record.category, fontWeight = FontWeight.Bold)
        Spacer(Modifier.weight(1f))
        Text(if (record.type == "income") "+${record.amount}" else "-${record.amount}",
             color = if (record.type == "income") Color(0xFF27AE60) else Color(0xFFC0392B))
    }
}
\end{lstlisting}

Zum Erfassen ist ein AlertDialog am bequemsten: ein paar Eingabefelder + Umschalter Einnahme/Ausgabe + Kategorie-Dropdown, und bei “OK” rufst du vm.addRecord(...) auf.

\begin{tcbwarn}{Achtung: Bei der Eingabeprüfung nicht faulenzen}
Das Betragsfeld musst du unbedingt prüfen: Leere Werte, Negativeingaben, “abc” und Ähnliches sind abzufangen (mit toDoubleOrNull() prüfen). Einer der häufigsten Gründe für App-Abstürze ist, “die Eingaben des Nutzers für brav zu halten”. Merke: Behandle alles, was Nutzer eintippen, wie etwas, das du feindlich prüfen musst.

\end{tcbwarn}

\booksection{Bearbeiten und Löschen: Daten müssen “rein und raus” können}

Ein Haushaltsbuch, das nur hinzufügen kann, ist nach drei Tagen unbrauchbar. Zwei Operationen noch:

\begin{lstlisting}[language=Kotlin]
// zwei Methoden im Dao ergänzen:
@Query("UPDATE records SET amount=:amount, category=:category, note=:note WHERE id=:id")
suspend fun update(id: Long, amount: Double, category: String, note: String)

@Query("DELETE FROM records WHERE id=:id")
suspend fun deleteById(id: Long)

// in der Oberfläche: langes Tippen auf einen Eintrag öffnet das Menü "Bearbeiten/Löschen" (in Compose: DropdownMenu oder AlertDialog)
// vor dem Löschen besser einen Bestätigungsdialog: "versehentliches Löschen" ist, was Nutzer von Ausgaben-Apps am meisten hassen
\end{lstlisting}

Zwei Erfahrungsregeln für die Bedienung: Beim Löschen immer eine zweite Bestätigung (gegen Versehen), beim Bearbeiten möglichst die Originalwerte voreintragen (ein Zeichen ändern ist leichter als alles neu tippen). Unterschätz diese Details nicht – “Produktgefühl” besteht genau aus solchen Details.

\booksection{Abnahme: drei Prüfebenen + zwei Fortgeschrittenen-Herausforderungen}

\booksection{Abnahme: drei Prüfebenen + zwei Fortgeschrittenen-Herausforderungen}

Wenn das Projekt fertig ist, nimm die Abnahme in drei Ebenen vor: “Funktion Grenzfälle Architektur”: Kannst du erfassen? Löschen? Sind die Daten nach einem Neustart noch da (Funktion)? Stürzt die App bei Betrag 0 oder negativ ab? Gehen bei einer Bildschirmdrehung Daten verloren (Grenzfälle)? Sind Oberfläche und Daten sauber getrennt, sind Felder hart kodiert (Architektur)?

Nach bestandener Basis-Abnahme gönn dir zwei “Fortgeschrittenen-Gerichte”:

\begin{enumerate}
  \item Monatsfilter: oben einen Monatswechsel ergänzen (z. B. 2026-08 v), die Liste zeigt nur den Verlauf des jeweiligen Monats;
  \item Kreisdiagramm nach Kategorien: mit Composes Canvas oder einer leichten Chart-Bibliothek "wohin das Geld fließt" als Tortendiagramm zeichnen – danach sieht das Haushaltsbuch schon richtig nach Produkt aus.
\end{enumerate}

\begin{tcbpractice}{Selbst ausprobieren: Diese Woche fertigstellen}
Arbeite in der Reihenfolge “Datenschicht ViewModel Listenseite Erfassungsseite Auswertungskarte”, jeden Tag ein Schritt. Wenn du festhängst, wirf die Fehlermeldung wörtlich in die Suchmaschine oder an eine KI (erinnerst du dich an Teil 7, “AI-Programmierung in der Praxis”? Beschreibe das Problem mit “erwartetes Verhalten, tatsächliches Verhalten, Schritte zur Reproduktion”). Wenn du fertig bist, führe eine Woche lang damit Buch – eine selbstgebaute App im eigenen Einsatz gibt dir Schwung für den weiteren Weg.

\end{tcbpractice}

\begin{tcbsummary}{Zusammenfassung des Kapitels}
\begin{itemize}
  \item Fünf Anforderungsfragen zerlegen in Schritte, von denen jeder einzeln funktioniert – dann loslegen
  \item Datenschicht: Record-Entity + Dao (inkl. SQL für Monatsstatistik) + Database
  \item MVVM-Arbeitsteilung: View zeichnet nur, ViewModel steuert die Logik, Model verwaltet die Speicherung
  \item Flow + stateIn + collectAsState: Daten ändern sich, die Oberfläche aktualisiert sich automatisch
  \item LazyColumn als träge ladende Liste; AlertDialog für die Erfassung
  \item Abnahme in drei Ebenen (Funktion/Grenzfälle/Architektur) + Zusatzaufgaben (Monatsfilter, Kreisdiagramm)
\end{itemize}

\end{tcbsummary}
